\chapter{Селектор направленного переноса в барионный сектор}
\label{ch:v8-baryon-directed-transfer-convention-selector-gate}

Предыдущая глава оставила два численных чтения барионного дискриминатора.
Здесь двусмысленность проверяется не сравнением десятичных значений, а
прямым ограничением канонического КМС-генератора на исходный угол $Q_L$.
Все вычисления выполняются в поле $\mathbb Q(x)$, где $x=e^{-2}$.

\section{Комплексное факторпространство стрелок}

Пусть $E_c:Q_L\to Y_R$, $c=1,2,3$, --- три фробениусово
нормированные цветовые стрелки,
\begin{equation}
 E_c=\frac1{\sqrt2}\langle c|\otimes I_2,
 \qquad
 \Tr(E_c^*E_d)=\delta_{cd}.
 \label{eq:v8-baryon-qlyr-normalized-arrows}
\end{equation}
Вещественный список $(E_c,-iE_c)$ содержит шесть элементов, но его
комплексный ранг равен трём. Для каждой пары грамов блок общего КМС-следа
и его псевдообратная матрица Мура--Пенроуза равны
\begin{equation}
 G_c=(a+b)
 \begin{pmatrix}1&-i\\ i&1\end{pmatrix},
 \qquad
 G_c^+=\frac1{4(a+b)}
 \begin{pmatrix}1&-i\\ i&1\end{pmatrix}.
 \label{eq:v8-baryon-qlyr-gram-pseudoinverse}
\end{equation}
Поэтому базис-независимый квадратичный оператор на исходном углу есть
\begin{equation}
 C_Q=\frac1{a+b}\sum_{c=1}^{3}E_c^*E_c
 =\frac{I_6}{2(a+b)}.
 \label{eq:v8-baryon-qlyr-source-casimir}
\end{equation}
Для связывающей инцидентности $A_0$ выполнены
$\|A_0\|_{\rm F}^2=13$ и
$(A_0^*A_0)|_{Q_L}=I_6$, откуда
\begin{equation}
 C_L=\frac{I_6}{13(a+b)}.
 \label{eq:v8-baryon-link-source-casimir}
\end{equation}

\section{КМС-направление и одночастичная пружина}

В направленном генераторе исходный угол умножается на скорость подъёма
$x$, а целевой --- на единичную скорость спуска. Следовательно,
\begin{equation}
 J_Q^{\rm src}=xC_Q=\frac{x}{2(a+b)}I_6,
 \qquad
 J_L^{\rm src}=xC_L=\frac{x}{13(a+b)}I_6.
 \label{eq:v8-baryon-directed-source-springs}
\end{equation}
Множитель $x$ является частью КМС-баланса и не может быть удалён
перенормировкой $1/x$ при сохранении того же направленного генератора.

\begin{proposition}[Селектор исходного угла]
Общий КМС-след и направленный генератор однозначно выбирают на $Q_L$
коэффициенты \eqref{eq:v8-baryon-directed-source-springs}. В частности,
чтение, удаляющее $x$, не является ограничением канонического генератора.
\end{proposition}

\section{Обнаруженный дефект половинного веса}

В проверяемом трёхкварковом отображении использовались пружины
\begin{equation}
 \widetilde J_Q=\frac{x}{4(a+b)}I_6,
 \qquad
 \widetilde J_L=\frac{x}{13(a+b)}I_6.
 \label{eq:v8-baryon-supplied-transfer-springs}
\end{equation}
Линкинговая часть совпадает с прямым ограничением, тогда как стрелочная
часть меньше ровно вдвое:
\begin{equation}
 \frac{J_Q^{\rm src}}{\widetilde J_Q}=2,
 \qquad
 \frac{J_L^{\rm src}}{\widetilde J_L}=1.
 \label{eq:v8-baryon-transfer-required-multipliers}
\end{equation}
Тем самым прежние два значения $0.2570$ и $0.1962$ не являются двумя
равноправными ограничениями одного и того же канонического оператора.
Первое сохраняет КМС-направление, но теряет половину стрелочного оператора;
второе дополнительно удаляет скорость $x$.

Если оставить неизменной уже проверенную калибровочную парную форму и
заменить только одночастичную переносную часть её прямым ограничением, то
\begin{align}
 r_1^{\rm dir}
 &=\frac{(10x+11)(375x^3+3916x^2+7267x+2782)}
 {26(x+1)(x+5)(25x+13)},
 \\
 v_{\rm dir}(x)
 &=\frac{52(25x^2+38x+13)}
 {375x^3+3916x^2+7267x+2782}.
 \label{eq:v8-baryon-direct-restriction-discriminator}
\end{align}
При $x=e^{-2}$ последнее выражение равно
$0.252005943190\ldots$; десятичная запись приведена только для чтения.

\begin{nogocriterion}[Граница селектора]
Гейт выводит каноническую одночастичную форму потерь на $Q_L$ и выявляет
точный дефект множителя два в прежнем стрелочном члене. Он ещё не выводит
полный функтор из одночастичного КМС-генератора в трёхчастичный массовый
оператор. Поэтому формула \eqref{eq:v8-baryon-direct-restriction-discriminator}
является каноническим кандидатом прямого ограничения, но не физической
теоремой о барионных массах.
\end{nogocriterion}

\begin{protocol}
\begin{enumerate}
 \item поле вычислений: \textbf{$\mathbb Q(x)$ без чисел с плавающей точкой};
 \item комплексный ранг шести вещественных стрелок: \textbf{$3$};
 \item псевдообратное грамово тождество: \textbf{точно};
 \item стрелочный оператор исходного угла: \textbf{$I_6/[2(a+b)]$};
 \item связывающий оператор исходного угла: \textbf{$I_6/[13(a+b)]$};
 \item сохранение КМС-множителя $x$: \textbf{обязательно};
 \item требуемый стрелочный множитель: \textbf{$2$};
 \item требуемый связывающий множитель: \textbf{$1$};
 \item прежнее базовое чтение: \textbf{не прямое каноническое ограничение};
 \item чтение с $1/x$: \textbf{не прямое каноническое ограничение};
 \item полный трёхчастичный подъём: \textbf{остаётся открыть};
 \item следующий гейт: \textbf{нормировка трёхчастичного подъёма,
 выполнена в следующей главе}.
\end{enumerate}
\end{protocol}

Машинный сертификат сохранён в
\path{s2t_v8_baryon_directed_transfer_convention_selector_gate_results.json}.